% =============================================================================
% 06_software.tex — lane WC
% Sources: tree layout, include/tensixflow/*.td, lib/tensixflow/**, tools/*,
%          kernels/handwritten_baselines/check_kernels_compile.sh,
%          python/tensixflow/**, test/, benchmarks/
% =============================================================================
\section{Software}\label{sec:software}

TensixFlow is a from-scratch C++/MLIR compiler stack: a custom
\texttt{tensixflow} dialect with tensor-level and fused ops, a
tile-level \emph{schedule IR}, seven lowering passes, a TT-Metalium-style
kernel emitter, a functional tile simulator, a static performance model with
an offline HTML visualizer, and a Python DSL plus a NumPy golden/correctness
harness. Everything builds and runs on a CPU-only box.

\subsection{Architecture}\label{sec:software-architecture}

Figure~\ref{fig:software-architecture} is the system map. The spine is
\texttt{tensixflow-opt}: hand-written MLIR (or MLIR text emitted by the
Python DSL) runs through five lowering passes, ending in the schedule IR.
That IR is the pivot: the emitter turns it into generated TT-Metalium-style
kernel projects (validated \emph{structurally only}), while the functional
tile simulator executes it directly (validated \emph{numerically} against an
fp32-exact NumPy golden). A side track exports the plan as JSON, runs the
static performance model, and renders the self-contained HTML report; a
compat track emits TTIR-shaped text.

\begin{figure}[htbp]
  \centering
  \resizebox{\textwidth}{!}{%
  \begin{tikzpicture}[
      box/.style={draw=tensixnavy, rounded corners=2pt, align=center,
                  font=\small, inner sep=3pt, fill=white},
      pass/.style={box, fill=tensixlight, minimum width=6.4cm, font=\small},
      outA/.style={box, fill=tensixteal!15, minimum width=4.4cm},
      outB/.style={box, fill=tensixamber!15, minimum width=4.6cm},
      side/.style={box, fill=tensixslate!10, minimum width=4.0cm},
      tag/.style={font=\scriptsize\itshape, tensixslate},
      arr/.style={-{Stealth[length=2.4mm]}, tensixnavy},
    ]
    % ---- inputs -------------------------------------------------------------
    \node[box, minimum width=5.2cm] (mlir) at (-3.6,0)
      {hand-written MLIR input\\[-1pt]
       \scriptsize\texttt{examples/rmsnorm\_matmul.mlir}, \texttt{mlp\_block.mlir}};
    \node[box, minimum width=5.6cm] (dsl) at (3.6,0)
      {Python DSL: \texttt{@tensixflow.jit} (thin builder)\\[-1pt]
       \scriptsize explicit builder calls --- emits MLIR text (no tracing)};

    % ---- pass spine ----------------------------------------------------------
    \node[pass, minimum width=5.0cm] (p1) at (0,-2.2) {\texttt{tensixflow-infer-tile-layout}};
    \node[pass, minimum width=5.0cm, below=0.26cm of p1] (p2)
      {\texttt{tensixflow-fuse-rmsnorm-matmul}\\[-1pt]
       \scriptsize (+ MLP chain $\Rightarrow$ \texttt{fused\_mlp})};
    \node[pass, minimum width=5.0cm, below=0.26cm of p2] (p3) {\texttt{tensixflow-plan-noc-sharding}};
    \node[pass, minimum width=5.0cm, below=0.26cm of p3] (p4) {\texttt{tensixflow-allocate-circular-buffers}};
    \node[pass, minimum width=5.0cm, below=0.26cm of p4] (p5) {\texttt{tensixflow-lower-to-ttkernel}};
    \begin{scope}[on background layer]
      \node[fit=(p1)(p5), draw=tensixnavy, densely dotted, fill=tensixnavy!4,
            inner sep=7pt] (pipebox) {};
    \end{scope}
    \node[tag, tensixnavy, font=\small\itshape, anchor=south] at (pipebox.north)
      {\texttt{tensixflow-opt} lowering spine};
    \draw[arr] (mlir.south) |- (p1.north west);
    \draw[arr] (dsl.south) |- (p1.north east);
    \node[tag] at (4.0,-1.25) {MLIR text};
    \draw[arr] (p1) -- (p2); \draw[arr] (p2) -- (p3);
    \draw[arr] (p3) -- (p4); \draw[arr] (p4) -- (p5);
    % L1 model contract feeding the sharding + CB passes
    \node[tag, align=center] (l1) at (-6.2,-4.75)
      {\textbf{L1 model contract}\\ \scriptsize \texttt{L1Model.h}: 1 MiB/core,\\
       \scriptsize 64 KiB reserve, 960 KiB CB/shard};
    \draw[arr, dashed] (l1.east) -- (p3.west);
    \draw[arr, dashed] (l1.east) -- (p4.west);

    % ---- schedule IR pivot ----------------------------------------------------
    \node[box, fill=tensixlight] (sched) at (0,-7.6)
      {\textbf{schedule IR} (\texttt{tensixflow.sched.*})\\[-1pt]
       \scriptsize \texttt{sched.kernel} $\supset$ reader / compute / writer tile programs:\\[-1pt]
       \scriptsize \texttt{tile\_loop} nests + CB + NoC + compute ops};
    \draw[arr] (p5) -- (sched) node[midway, right=2pt, tag] {output};

    % ---- track A: structural (left) -------------------------------------------
    \node[outA] (em) at (-5.3,-9.3)
      {\texttt{--tensixflow-emit-metalium}\\[-1pt]
       \scriptsize (\texttt{output-dir=<dir>})};
    \node[outA, below=0.45cm of em] (proj)
      {generated standalone project\\[-1pt]
       \scriptsize host/reader/compute/writer.cpp\\[-1pt]
       \scriptsize + CMakeLists.txt + README.md};
    \node[outA, below=0.45cm of proj] (chk)
      {compile-only stub check @\texttt{df36cd6}\\[-1pt]
       \scriptsize \texttt{check-generated-kernels-compile}\\[-1pt]
       \scriptsize \textbf{[STRUCTURAL ONLY]}};
    \draw[arr] (sched.west) -| (em.north) node[pos=0.25, above=1pt, tag] {emits};
    \draw[arr] (em) -- (proj);
    \draw[arr] (proj) -- (chk);
    \node[tag, align=center] at (-5.3,-13.15)
      {stub contract: signature-only headers,\\ every body empty, pinned @\texttt{df36cd6}};

    % ---- track B: numerical (center) -------------------------------------------
    \node[outB] (sim) at (0,-9.6)
      {functional tile simulator\\[-1pt]
       \scriptsize (\texttt{tensixflow-simulate})\\[-1pt]
       \scriptsize executes the schedule IR};
    \node[outB, below=0.45cm of sim] (gold)
      {harness vs fp32-exact\\[-1pt] NumPy golden\\[-1pt]
       \scriptsize PCC $\ge 0.9875$ --- \textbf{[NUMERICAL VALIDATION]}};
    \draw[arr] (sched) -- (sim);
    \draw[arr] (sim) -- (gold);

    % ---- side track: translate -> plan.json -> perf -> HTML (right) -------------
    \node[side, right=1.9cm of p4] (trans)
      {\texttt{tensixflow-translate}\\[-1pt]
       \scriptsize \texttt{--tensixflow-to-json}\\[-1pt]
       \scriptsize (IR stops pre-\texttt{lower-to-ttkernel})};
    \node[side, below=0.42cm of trans] (plan)
      {\texttt{plan.json}\\[-1pt]
       \scriptsize plan/perf/viz JSON schema\\[-1pt]
       \scriptsize (\texttt{schema/plan.schema.json})};
    \node[side, below=0.42cm of plan] (perf)
      {static perf model\\[-1pt]
       \scriptsize (\texttt{tensixflow-report} runs it\\[-1pt] before rendering)};
    \node[side, below=0.42cm of perf] (html)
      {offline HTML report\\[-1pt]
       \scriptsize grid, shards, NoC routes,\\[-1pt]
       \scriptsize CB usage, timeline, bottlenecks};
    \draw[arr] (p4.east) -- (trans.west);
    \draw[arr] (trans) -- (plan); \draw[arr] (plan) -- (perf); \draw[arr] (perf) -- (html);

    % ---- compat track ------------------------------------------------------------
    \node[side, right=1.9cm of p2] (ttir)
      {compat: \texttt{--tensixflow-emit-ttir}\\[-1pt]
       \scriptsize TTIR text only (\texttt{ttir.*}); unmapped\\[-1pt]
       \scriptsize $\Rightarrow$ documented UNSUPPORTED marker};
    \draw[arr, dashed] (p2.east) -- (ttir.west);

    % ---- contracts legend ----------------------------------------------------------
    \node[draw=tensixslate, rounded corners=2pt, align=left, font=\scriptsize,
          fill=tensixlight, inner sep=5pt] at (0,-14.6)
      {\textbf{Data contracts connecting the boxes:}\\
       \textbf{(1)} schedule IR (\texttt{TensixflowSchedule.td}) --- the same IR the emitter prints and the simulator executes;\\
       \textbf{(2)} plan/perf/viz JSON schema (\texttt{schema/plan.schema.json});\\
       \textbf{(3)} signature-only tt-metal stub headers pinned @\texttt{df36cd6};\\
       \textbf{(4)} L1 model constants (\texttt{L1Model.h}) shared by sharding, CB allocation, and the perf model.};
  \end{tikzpicture}}
  \caption{TensixFlow architecture: MLIR input and the Python DSL feed the
  \texttt{tensixflow-opt} lowering spine; the schedule IR is consumed twice ---
  by the Metalium emitter (structural validation) and by the functional tile
  simulator (numerical validation). The side track exports \texttt{plan.json},
  runs the perf model, and renders the offline HTML report; the compat track
  emits TTIR-only text.}
  \label{fig:software-architecture}
\end{figure}

\subsection{The \texttt{tensixflow} dialect: ops}

All ops live in \texttt{include/tensixflow/TensixflowOps.td}
(Table~\ref{tab:software-ops}); the schedule IR's \texttt{sched.*} ops are
defined separately in \texttt{TensixflowSchedule.td} and enumerated in the
Figure~\ref{fig:software-architecture} caption above.

\begin{table}[htbp]
  \centering
  \caption{Ops of the \texttt{tensixflow} dialect (tensor level), with their
  verifier / shape-inference status exactly as declared in
  \texttt{TensixflowOps.td}.}
  \label{tab:software-ops}
  \small
  \begin{tabular}{@{}p{4.3cm}p{6.2cm}p{4.6cm}@{}}
    \toprule
    \textbf{Op} & \textbf{Semantics} & \textbf{Verifier / shape inference} \\
    \midrule
    \texttt{func} & Function container; isolated from above, single block, implicit \texttt{return} terminator & FunctionOpInterface, Symbol; custom assembly \\
    \texttt{return} & Terminates \texttt{func}; Pure, ReturnLike & --- \\
    \texttt{constant} & SSA value from a typed attribute & \texttt{hasFolder} (constant folding) \\
    \texttt{add} & Elementwise sum of same-shape tensors & SameOperandsAndResultType; \texttt{hasCanonicalizer} (identity-add removal) \\
    \texttt{gelu} & Elementwise GELU/SiLU activation & SameOperandsAndResultType \\
    \texttt{rmsnorm} & $y = x \cdot \mathrm{rsqrt}(\mathrm{mean}(x^2)+\epsilon) \cdot \gamma$ over the last dimension & \texttt{hasVerifier}; \texttt{inferReturnTypeComponents} \\
    \texttt{matmul} & $[...,M,K] \times [K,N] \to [...,M,N]$ & \texttt{hasVerifier}; \texttt{inferReturnTypeComponents} \\
    \texttt{fused\_rmsnorm\_matmul} & $\mathrm{matmul}(\mathrm{rmsnorm}(x,\gamma), w)$ as one kernel & \texttt{hasVerifier}; \texttt{inferReturnTypeComponents} \\
    \texttt{fused\_mlp} & $(\mathrm{gelu}(\mathrm{matmul}(\mathrm{rmsnorm}(x,\gamma),w_1)+b_1)\, w_2) + \mathrm{residual}$ as one kernel & \texttt{hasVerifier}; \texttt{inferReturnTypeComponents} \\
    \bottomrule
  \end{tabular}
\end{table}

\subsection{Passes}

Registered by \texttt{tensixflow-opt} (Table~\ref{tab:software-passes}); the
CLI chain used end-to-end is the five-pass spine of
Figure~\ref{fig:software-architecture} followed by
\texttt{--tensixflow-emit-metalium=output-dir=<dir>}.

\begin{table}[htbp]
  \centering
  \caption{Passes of the TensixFlow transform pipeline, one line each, as
  declared in \texttt{TensixflowPasses.td}.}
  \label{tab:software-passes}
  \small
  \begin{tabular}{@{}p{5.4cm}p{9.7cm}@{}}
    \toprule
    \textbf{Pass flag} & \textbf{Purpose} \\
    \midrule
    \texttt{--tensixflow-infer-tile-layout} & Attach a \texttt{\#tensixflow.tile\_layout} encoding to each tensor result (dims rounded to $32\times32$ tiles, default $8\times8$ grid). \\
    \texttt{--tensixflow-canonicalize} & Canonicalize tensixflow ops: identity-add removal, trivial constant folding, structural dedup. \\
    \texttt{--tensixflow-fuse-rmsnorm-matmul} & Rewrite an adjacent rmsnorm$\to$matmul pair into one fused op, cutting intermediate DRAM traffic; the same pass grows the MLP chain into \texttt{fused\_mlp}. \\
    \texttt{--tensixflow-plan-noc-sharding} & Assign tensor shards to cores of the $8\times8$ grid under the L1 budget and plan NoC movement (consumes \texttt{L1Model.h} constants). \\
    \texttt{--tensixflow-allocate-circular-buffers} & Allocate L1 circular buffers for the tile streams between reader/compute/writer kernels. \\
    \texttt{--tensixflow-lower-to-ttkernel} & Lower fused/sharded ops into the schedule IR (tile-loop nests + CB + NoC + compute ops). \\
    \texttt{--tensixflow-emit-metalium} & Emit TT-Metalium-style host/reader/compute/writer \texttt{.cpp}; with \texttt{output-dir=} write a complete standalone project (+ CMakeLists.txt + README.md). \\
    \texttt{--tensixflow-emit-ttir} (compat) & Emit TTIR-shaped MLIR text to stdout --- \texttt{ttir.*} only (TTNN/D2M stay prose mappings); unmapped ops emit a documented UNSUPPORTED marker, transitively. \\
    \bottomrule
  \end{tabular}
\end{table}

The lit suite also carries one TTIR parse test that feeds the emitted text to
a real \texttt{ttmlir-opt}; since tt-mlir's toolchain is not built on this
box, that test is gated and reports UNSUPPORTED
(Table~\ref{tab:software-validation-tiers}, tier 1).

\subsection{Validation tiers (live counts)}\label{sec:software-validation}

All counts were re-run on this box for this report and reproduce the F3
final-wave results (Appendix, Table~\ref{tab:appendix-f3-commands}).

\begin{table}[htbp]
  \centering
  \caption{Validation tiers with current counts (verified live 2026-08-03).
  C/C++ tiers run under \texttt{nix develop --command bash -c '...'} so the
  pinned LLVM/MLIR 20.1.8 environment is on the PATH.}
  \label{tab:software-validation-tiers}
  \small
  \begin{tabular}{@{}p{2.6cm}p{7.4cm}p{5.1cm}@{}}
    \toprule
    \textbf{Tier} & \textbf{Command} & \textbf{Current result} \\
    \midrule
    1. lit/FileCheck IR suite & \texttt{cmake --build build --target check-tensixflow} & Total 25: \textbf{24 passed + 1 unsupported} (\texttt{emit\_ttir\_ttmlir\_parse.mlir}, gated on a \texttt{ttmlir-opt} this box does not build) \\
    2. ctest C++ units (gtest) & \texttt{cd build \&\& ctest --output-on-failure} & \textbf{11/11 passed}: plan\_schema, scheduler\_model, grid\_planner, weight\_planner, cb\_allocator, noc\_estimator, plan\_report, rounding, tile\_simulator, perf\_model, hw-acceptance \\
    3. pytest (Python) & \texttt{uv run --with numpy --with pytest pytest python/tensixflow/tests/ -q} & \textbf{89 passed} (golden, harness, DSL, lowering, rounding, correctness) \\
    4. generated-kernel compile check & \texttt{cmake --build build --target check-generated-kernels-compile} & baseline compiles, legacy \texttt{EnqueueProgram} fixture rejected, 4 generated kernels compile (also runs inside ctest's \texttt{hw-acceptance}) \\
    \bottomrule
  \end{tabular}
\end{table}

Two notes on tier 3. The harness runs under \texttt{uv}'s ephemeral
environment because no system Python ships \texttt{numpy} and \texttt{pytest}
together here. And there is a documented ordering dependency: the rounding
gtest \emph{writes} a file the rounding pytest \emph{reads}, so
\texttt{ctest -R rounding} must run before the pytest tier (as the ctest
target ordering above already does). Tier 3 includes the
simulator-vs-golden correctness harness (\texttt{test\_correctness.py}, 11
tests): both targets clear the PCC floor of 0.9875 with non-vacuity defenses
(maximum absolute difference strictly positive on every run, fp32-golden vs
bf16-simulator structural assertion, and per-target empirical trip-wires).

\subsection{Generated-kernel validation mechanics}\label{sec:software-stubs}

The generated TT-Metalium kernels are never executed and are validated
\emph{structurally only}, by compile-checking them against stub headers:

\begin{itemize}[nosep]
  \item \textbf{Signature-only stubs, pinned @\texttt{df36cd6}.}
        \texttt{kernels/handwritten\_baselines/tt\_metal\_stubs/} mirrors the
        tt-metal kernel- and host-side API surface the generated kernels use
        (\texttt{dataflow\_api.h}, \texttt{compute\_kernel\_api.h} and its
        \texttt{matmul.h}/\texttt{pack.h}/\texttt{reg\_api.h} sub-includes,
        \texttt{host\_api.hpp}, \texttt{kernel\_types.hpp}). Every body is
        empty (\texttt{\{\}} / \texttt{return \{\}}) --- a compile contract
        only, no behavior; the headers say so verbatim.
  \item \textbf{The check itself} is orchestrated by
        \texttt{kernels/handwritten\_baselines/check\_kernels\_compile.sh} as
        the CMake target \texttt{check-generated-kernels-compile}: compile
        each translation unit with
        \texttt{g++ -std=c++20 -fsyntax-only -Wall -I <stubs>}.
        Syntax-only means no link step and no device is possible --- but the
        compiler still resolves every API name and signature the kernels use.
  \item \textbf{Non-vacuous by construction.} The target asserts three
        things, all required: (1) \emph{positive} --- the current-API
        baseline kernel compiles, proving the stubs still carry the current
        symbols; (2) \emph{negative} --- a fixture referencing the removed
        legacy symbol \texttt{EnqueueProgram} must \emph{fail} with a
        diagnostic naming that symbol, proving the stubs have not drifted
        back onto the legacy API; (3) every generated kernel file compiles.
        If the check ever passed vacuously, assertion (2) would fail instead.
\end{itemize}
